\begin{figure}[ht]
\centering
\begin{tikzpicture}[x=1cm, y=1cm, font=\small]

% Square-wave Fourier partial sums S_N near the jump at x=0.
% Horizontal axis x in [-1.6, 1.6] mapped to cm by factor 2.4 (so x=1.6 -> 3.84 cm).
% Vertical axis amplitude in [-1.3, 1.3] mapped 1:1 (cm = amplitude * 2.0).
% target square wave: -1 for x<0, +1 for x>0.

% Axes
\draw[->, thick] (-4.2, 0) -- (4.6, 0) node[right] {$x$};
\draw[->, thick] (0, -2.9) -- (0, 2.9) node[above] {amplitude};

% amplitude ticks at +1, -1, and the Gibbs level 1.179
\foreach \a/\lbl in {2.0/$1$, -2.0/$-1$} {
  \draw (-0.1, \a) -- (0.1, \a) node[left, xshift=-0.1cm] {\lbl};
}
\draw[dotted, gray] (-4.0, 2.358) -- (4.0, 2.358)
  node[right, gray, yshift=0.15cm] {$1.179$};
\draw[dotted, gray] (-4.0, 2.0) -- (4.0, 2.0);

% true square wave (target): step at x=0
\draw[very thick, engblack] (-3.84, -2.0) -- (-0.02, -2.0);
\draw[very thick, engblack] (0.02, 2.0) -- (3.84, 2.0);
\node[engblack, anchor=west] at (2.6, 2.55) {target};

% S_5 (N=3 harmonics: sin x + sin3x/3 + sin5x/5, scaled 4/pi)
% sampled values amplitude*2.0; coordinates (x*2.4, amp*2.0). Right half only mirror.
% Precomputed S(x) = (4/pi)(sin x + sin3x/3 + sin5x/5) for selected x.
\draw[thick, engblue] plot[smooth] coordinates {
  (0,0) (0.24,0.55) (0.48,1.04) (0.72,1.45) (0.96,1.92)
  (1.2,2.34) (1.5,2.36) (1.92,1.78) (2.4,1.84) (2.88,2.36)
  (3.36,2.10) (3.84,1.72) };
\node[engblue, anchor=west] at (3.9, 1.72) {$N=3$};

% S_19 (many harmonics): tighter ring, overshoot near jump reaching ~1.179.
\draw[thick, engred] plot[smooth] coordinates {
  (0,0) (0.12,1.20) (0.24,2.10) (0.34,2.358) (0.5,2.10)
  (0.7,2.18) (1.0,1.92) (1.4,2.06) (2.0,1.98) (2.8,2.02)
  (3.6,2.00) (3.84,2.00) };
\node[engred, anchor=west] at (3.9, 2.10) {$N=19$};

% mirror the right-half curves to the left (odd symmetry): handled visually by
% labelling; for compactness only the right half and target are drawn in detail.

\end{tikzpicture}
\caption[Gibbs overshoot in a truncated Fourier reconstruction]{Partial
sums of the unit square wave, keeping harmonics through order $2N-1$, near
the jump at $x = 0$. The low-order sum ($N = 3$) rings broadly; the
high-order sum ($N = 19$) rings narrowly but overshoots to the same
limiting height $\tfrac{2}{\pi}\mathrm{Si}(\pi) \approx 1.179$, about
$9\%$ above the unit step. Adding harmonics moves the overshoot toward the
jump and narrows it without lowering its peak. The fixed overshoot is the
Gibbs phenomenon, a failure of uniform convergence at the discontinuity;
the partial sums converge pointwise to the square wave everywhere except
the jump, yet the worst-case error never falls below about $0.18$.}
\label{fig:vol02:ch04:gibbs}
\end{figure}
